% GENERATED FILE. Edit canonical lessons, then run npm run generate:books.
% canonical-source-hash: fnv1a64:f342fd804a0e0480
% canonical-lessons: AR-C315-nasr, AR-C315-hazima, AR-C315-intasara, AR-C315-hazama-defeat, AR-C315-istaslama

\chapter{Victory and defeat}
\label{ch:ar-victory-and-defeat}
\hlchaptermodality{315}

\begin{chapteropening}
\textbf{By the end of this chapter:} \emph{I can say a victory in battle, a defeat, to triumph, to defeat and to surrender.}

Complete the last lesson of chapter 315: I can say a victory in battle, a defeat, to triumph, to defeat and to surrender.
\end{chapteropening}

\section[\texorpdfstring{\ar{نصر} (naṣr)}{نصر (naṣr)}]{\ar{نصر} (naṣr) — a victory in battle}

\label{lesson:AR-C315-nasr}



\begin{quote}
Before the new one: say the Arabic for to escort a coffin, then the Arabic for a grave.
\end{quote}

\subsection*{You'll want to know: \ar{نصر}}
\textbf{\ar{نصر}} — \emph{naṣr} — \textquotedblleft{}a victory in battle\textquotedblright{}.

\begin{quote}
\ar{احتفل} \ar{الناس} \ar{بالنصر} \ar{في} \ar{الشوارع}.
\end{quote}

\begin{quote}
\emph{iḥtafala an-nās bi-n-naṣr fī ash-shawāriʿ.}
\end{quote}

\begin{quote}
People celebrated the victory in the streets.
\end{quote}

\begin{cousinweb}
\textbf{\ar{نصر}} (\emph{naṣr}) is built on the native Arabic root \emph{n-ṣ-r} (Lane), \textquotedblleft{}to aid someone against an enemy\textquotedblright{}. Help that wins the day is victory.
\end{cousinweb}

\begin{grammarlens}[title={What you've built}]
One more word for telling what happened.
\end{grammarlens}

\subsection*{Your turn}
\begin{itemize}
\raggedright
  \item \emph{Say it:} \emph{naṣr}
  \item \emph{Say it:} \emph{naṣr}, once more
  \item \emph{Say it:} \emph{iḥtafala an-nās bi-n-naṣr fī ash-shawāriʿ.} — People celebrated the victory in the streets.
  \item \emph{Recall:} say the Arabic for to escort a coffin, then the Arabic for a grave, then say \emph{naṣr} again
\end{itemize}

\begin{tcolorbox}[breakable,colback=teal!4,colframe=teal!35!black,title={Before you move on}]
What does \ar{نصر} mean? (\textquotedblleft{}A victory in battle\textquotedblright{}.) Say it once more.
\end{tcolorbox}
\section[\texorpdfstring{\ar{هزيمة} (hazīma)}{هزيمة (hazīma)}]{\ar{هزيمة} (hazīma) — a defeat}

\label{lesson:AR-C315-hazima}



\begin{quote}
Before the new one: say the Arabic for a grave, then the Arabic for a victory in battle.
\end{quote}

\subsection*{You'll want to know: \ar{هزيمة}}
\textbf{\ar{هزيمة}} — \emph{hazīma} — \textquotedblleft{}a defeat\textquotedblright{}.

\begin{quote}
\ar{كانت} \ar{الهزيمة} \ar{صعبة} \ar{على} \ar{الفريق}.
\end{quote}

\begin{quote}
\emph{kānat al-hazīma ṣaʿba ʿalā al-farīq.}
\end{quote}

\begin{quote}
The defeat was hard on the team.
\end{quote}

\begin{cousinweb}
\textbf{\ar{هزيمة}} (\emph{hazīma}) is built on the native Arabic root \emph{h-z-m} (Lane), \textquotedblleft{}to rout, to put to flight\textquotedblright{}. \emph{Hazīma} is the rout of the side that loses.
\end{cousinweb}

\begin{grammarlens}[title={What you've built}]
One more word for telling what happened.
\end{grammarlens}

\subsection*{Your turn}
\begin{itemize}
\raggedright
  \item \emph{Say it:} \emph{hazīma}
  \item \emph{Say it:} \emph{hazīma}, once more
  \item \emph{Say it:} \emph{kānat al-hazīma ṣaʿba ʿalā al-farīq.} — The defeat was hard on the team.
  \item \emph{Recall:} say the Arabic for a grave, then the Arabic for a victory in battle, then say \emph{hazīma} again
\end{itemize}

\begin{tcolorbox}[breakable,colback=teal!4,colframe=teal!35!black,title={Before you move on}]
What does \ar{هزيمة} mean? (\textquotedblleft{}A defeat\textquotedblright{}.) Say it once more.
\end{tcolorbox}
\section[\texorpdfstring{\ar{انتصر} (intaṣara)}{انتصر (intaṣara)}]{\ar{انتصر} (intaṣara) — to win a victory; to triumph}

\label{lesson:AR-C315-intasara}



\begin{quote}
Before the new one: say the Arabic for a victory in battle, then the Arabic for a defeat.
\end{quote}

\subsection*{You'll want to know: \ar{انتصر}}
\textbf{\ar{انتصر}} — \emph{intaṣara} — \textquotedblleft{}to win a victory; to triumph\textquotedblright{}.

\begin{quote}
\ar{انتصر} \ar{فريقنا} \ar{في} \ar{المباراة} \ar{الأخيرة}.
\end{quote}

\begin{quote}
\emph{intaṣara farīqunā fī al-mubārāh al-akhīra.}
\end{quote}

\begin{quote}
Our team won the last match.
\end{quote}

Present: \ar{ينتصر}, \emph{yantaṣir}.

\ar{انتصر} \ar{على}: to triumph over.

\begin{cousinweb}
\textbf{\ar{انتصر}} (\emph{intaṣara}) is built on the native Arabic root \emph{n-ṣ-r} (Lane), \textquotedblleft{}to help against an enemy\textquotedblright{}. It is Form VIII, \textquotedblleft{}to get the victory\textquotedblright{}.
\end{cousinweb}

\begin{grammarlens}[title={What you've built}]
One more word for telling what happened.
\end{grammarlens}

\subsection*{Your turn}
\begin{itemize}
\raggedright
  \item \emph{Say it:} \emph{intaṣara}
  \item \emph{Say it:} \emph{intaṣara}, once more
  \item \emph{Say it:} \emph{intaṣara farīqunā fī al-mubārāh al-akhīra.} — Our team won the last match.
  \item \emph{Recall:} say the Arabic for a victory in battle, then the Arabic for a defeat, then say \emph{intaṣara} again
\end{itemize}

\begin{tcolorbox}[breakable,colback=teal!4,colframe=teal!35!black,title={Before you move on}]
What does \ar{انتصر} mean? (\textquotedblleft{}To win a victory; to triumph\textquotedblright{}.) Say it once more.
\end{tcolorbox}
\section[\texorpdfstring{\ar{هزم} (hazama)}{هزم (hazama)}]{\ar{هزم} (hazama) — to defeat}

\label{lesson:AR-C315-hazama-defeat}



\begin{quote}
Before the new one: say the Arabic for a defeat, then the Arabic for to win a victory.
\end{quote}

\subsection*{You'll want to know: \ar{هزم}}
\textbf{\ar{هزم}} — \emph{hazama} — \textquotedblleft{}to defeat\textquotedblright{}.

\begin{quote}
\ar{هزم} \ar{الفريق} \ar{الصغير} \ar{الفريق} \ar{الكبير}.
\end{quote}

\begin{quote}
\emph{hazama al-farīq aṣ-ṣaghīr al-farīq al-kabīr.}
\end{quote}

\begin{quote}
The small team beat the big team.
\end{quote}

Present: \ar{يهزم}, \emph{yahzim}.

The passive \ar{هزم}, \emph{huzima}, is to be defeated.

\begin{cousinweb}
\textbf{\ar{هزم}} (\emph{hazama}) is built on the native Arabic root \emph{h-z-m} (Lane), \textquotedblleft{}to rout; to put an army to flight\textquotedblright{}: Lane gives \emph{inhazama}, of an army that was routed.
\end{cousinweb}

\begin{grammarlens}[title={What you've built}]
One more word for telling what happened.
\end{grammarlens}

\subsection*{Your turn}
\begin{itemize}
\raggedright
  \item \emph{Say it:} \emph{hazama}
  \item \emph{Say it:} \emph{hazama}, once more
  \item \emph{Say it:} \emph{hazama al-farīq aṣ-ṣaghīr al-farīq al-kabīr.} — The small team beat the big team.
  \item \emph{Recall:} say the Arabic for a defeat, then the Arabic for to win a victory, then say \emph{hazama} again
\end{itemize}

\begin{tcolorbox}[breakable,colback=teal!4,colframe=teal!35!black,title={Before you move on}]
What does \ar{هزم} mean? (\textquotedblleft{}To defeat\textquotedblright{}.) Say it once more.
\end{tcolorbox}
\section[\texorpdfstring{\ar{استسلم} (istaslama)}{استسلم (istaslama)}]{\ar{استسلم} (istaslama) — to surrender, to give in}

\label{lesson:AR-C315-istaslama}



\begin{quote}
Before the new one: say the Arabic for to win a victory, then the Arabic for to defeat.
\end{quote}

\subsection*{You'll want to know: \ar{استسلم}}
\textbf{\ar{استسلم}} — \emph{istaslama} — \textquotedblleft{}to surrender, to give in\textquotedblright{}.

\begin{quote}
\ar{استسلم} \ar{الجنود} \ar{بعد} \ar{يومين}.
\end{quote}

\begin{quote}
\emph{istaslama al-junūd baʿda yawmayn.}
\end{quote}

\begin{quote}
The soldiers surrendered after two days.
\end{quote}

Present: \ar{يستسلم}, \emph{yastaslim}.

\ar{لا} \ar{تستسلم}! Do not give up!

\begin{cousinweb}
\textbf{\ar{استسلم}} (\emph{istaslama}) is built on the native Arabic root \emph{s-l-m} (Lane), \textquotedblleft{}to be safe, sound\textquotedblright{}, the root of \emph{salām}, peace. It is Form X, \textquotedblleft{}to yield oneself up\textquotedblright{}.
\end{cousinweb}

\begin{grammarlens}[title={What you've built}]
One more word for telling what happened.
\end{grammarlens}

\subsection*{Your turn}
\begin{itemize}
\raggedright
  \item \emph{Say it:} \emph{istaslama}
  \item \emph{Say it:} \emph{istaslama}, once more
  \item \emph{Say it:} \emph{istaslama al-junūd baʿda yawmayn.} — The soldiers surrendered after two days.
  \item \emph{Recall:} say the Arabic for to win a victory, then the Arabic for to defeat, then say \emph{istaslama} again
\end{itemize}

\begin{tcolorbox}[breakable,colback=teal!4,colframe=teal!35!black,title={Before you move on}]
What does \ar{استسلم} mean? (\textquotedblleft{}To surrender, to give in\textquotedblright{}.) Say it once more.
\end{tcolorbox}
